\subsection*{a.3 -- Optimized dimensions}

\paragraph{Formulation.} Maximize the achievable peak sensitivity,
Eq.~\ref{eq:Smax}, $|H|_{\mathrm{max}}\propto L\,e^{-bg_{0}}/\VDCMAX$,
over all eleven geometric parameters, subject to
\begin{itemize}
  \item the four assignment constraints: $f_{r}>\SI{500}{\kilo\hertz}$,
        $\VPI>\SI{200}{\volt}$, $\VDCMAX<\VPI$,
        $a_{0,\mathrm{max}}(\omega_{r})>\SI{2e4}{\meter\per\second\squared}$;
  \item the travel condition Eq.~\ref{eq:travel} with 20\% margin;
  \item the four side-stability channels (a.2) with 15\% voltage margin over
        $\VDCMAX$;
  \item internal-mode separation: bar/spine tip, finger and beam bending
        modes $\ge 2.5f_{r}$ (so that the first eigenmode is cleanly the
        actuation mode and the rigid-shuttle assumption holds);
  \item the comb fitting under the spine, $\approx N(W_{c}+h)+6W_{c}\le L$;
  \item technology: all features in
        [\SI{200}{\nano\meter}, \SI{1}{\milli\meter}], $2<N<200$,
        $t=\SI{200}{\nano\meter}$ fixed.
\end{itemize}
The model is the Part-1 lumped element set with figure-faithful masses
(validated against the Part-2 FEM in Q1/Q2 there), $Q$ from the Part-1
damping model ($b_{\mathrm{air}}=c_{\mathrm{air}}A_{m}$,
$Q_{\mathrm{clamp}}=500$), and a full grid search over $\sim10^{6}$
candidates followed by rounding to clean values
(\texttt{analysis/p3a\_optimize.py}).

\paragraph{Design logic.} The search converges onto a clear logic:
\begin{enumerate}
  \item $g_{0}$ stays at the \SI{200}{\nano\meter} minimum (sensitivity
        $\propto e^{-bg_{0}}$);
  \item the suspension is softened exactly to the resonance floor --
        $\VDCMAX\propto\sqrt{kh/N}$, and $k$ may drop as far as
        $f_{r}>\SI{500}{\kilo\hertz}$ allows. Conveniently, at
        $g_{0}=\SI{200}{\nano\meter}$ the acceleration constraint
        $a_{0,\mathrm{max}}(\omega_{r})=g_{0}\omega_{r}^{2}/Q$ with
        $Q\approx\omega_{r}\rho t/c_{\mathrm{air}}\approx80$ is satisfied at
        almost exactly the same frequency, so both floors coincide;
  \item the comb is completely re-proportioned: short, stubby fingers
        ($L_{c}\times W_{c}=5.5\times\SI{0.7}{\micro\meter}$, internal mode
        \SI{32}{\mega\hertz}) at a small rest overlap $d_{0}=\SI{1}{\micro\meter}$
        leave a tip clearance $g_{t0}=\SI{4.5}{\micro\meter}$ and a stable
        travel $x_{\mathrm{PI}}=\SI{3.07}{\micro\meter}$, comfortably beyond
        the \SI{1.67}{\micro\meter} optimum; $N=58$ gaps at
        $h=\SI{0.9}{\micro\meter}$ fill the spine length;
  \item the plates slim down to $W_{m}=\SI{10}{\micro\meter}$ -- the
        spine-bending (negative elastic foundation) limit; thinner plates
        would buckle sideways electrostatically far below the operating
        voltage;
  \item $L$ stays at \SI{100}{\micro\meter}: the plate mass required for
        spine stability grows $\propto L^{7/3}$, so the extra optical length
        of a longer bar is eaten by the higher $\VDCMAX$ -- the objective is
        nearly flat in $L$, and the shortest stiff device wins on robustness.
\end{enumerate}

\begin{figure}[H]
\centering
\includegraphics[width=0.8\textwidth]{figures/p3a_design_space.pdf}
\caption{The $\sim$47\,000 feasible designs of the refined search in the
sensitivity--$\VDCMAX$ plane (color: modulation length $L$). Each arc is one
$L$-family; the objective $|H|_{\mathrm{max}}\propto L/\VDCMAX$ improves
toward the top left. The chosen design (star) sits on the frontier; the
$\approx$1\% gap to the best raw grid point is the cost of rounding to clean
dimensions.}
\label{fig:designspace}
\end{figure}

\begin{table}[H]
\centering
\small
\caption{Final dimensions (the ``Final Value'' column of the assignment
table) next to the Part-1/2 values.}
\label{tab:dims}
\begin{tabular}{llcc}
\toprule
Feature & Symbol & Part 1/2 & \textbf{Part 3a final}\\
\midrule
Length of modulation section & $L$       & \SI{100}{\micro\meter} & \textbf{\SI{100}{\micro\meter}}\\
Modulation gap at rest       & $g_{0}$   & \SI{200}{\nano\meter}  & \textbf{\SI{200}{\nano\meter}}\\
Length of beams              & $L_{b}$   & \SI{75}{\micro\meter}  & \textbf{\SI{40}{\micro\meter}}\\
Width of beams               & $W_{b}$   & \SI{5}{\micro\meter}   & \textbf{\SI{1.85}{\micro\meter}}\\
Length of mass central section & $L_{m}$ & \SI{40}{\micro\meter}  & \textbf{\SI{30}{\micro\meter}}\\
Width of mass sections       & $W_{m}$   & \SI{30}{\micro\meter}  & \textbf{\SI{10}{\micro\meter}}\\
Comb-drive gap               & $h$       & \SI{1}{\micro\meter}   & \textbf{\SI{0.9}{\micro\meter}}\\
Comb-drive overlap at rest   & $d_{0}$   & \SI{19}{\micro\meter}  & \textbf{\SI{1.0}{\micro\meter}}\\
Comb-drive finger length     & $L_{c}$   & \SI{20}{\micro\meter}  & \textbf{\SI{5.5}{\micro\meter}}\\
Comb-drive finger width      & $W_{c}$   & \SI{1}{\micro\meter}   & \textbf{\SI{0.7}{\micro\meter}}\\
Bottom end fingers           & $3W_{c}$  & \SI{3}{\micro\meter}   & \textbf{\SI{2.1}{\micro\meter}}\\
Thickness of suspended parts & $t$       & \SI{200}{\nano\meter}  & \textbf{\SI{200}{\nano\meter}} (fixed)\\
Number of comb gaps          & $N$       & 40                     & \textbf{58}\\
\bottomrule
\end{tabular}
\end{table}

\begin{table}[H]
\centering
\small
\caption{Predicted performance (analytical model) and constraint check.}
\label{tab:perf}
\begin{tabular}{lccc}
\toprule
quantity & Part 1/2 & Part 3a & constraint\\
\midrule
$k$ & \SI{40.3}{\newton\per\meter} & \SI{13.5}{\newton\per\meter} & --\\
$\meff$ & \SI{3.80}{\pico\gram} & \SI{1.18}{\pico\gram} & --\\
$f_{r}$ & 518 kHz (FEM 496) & 539 kHz (FEM 522) & $>\SI{500}{\kilo\hertz}$ \checkmark\\
$\VDCMAX$ & \SI{1377}{\volt} & \textbf{\SI{627}{\volt}} & $<\VPI$ \checkmark (margin 1.15)\\
$\VPI$ axial / side & 347 / \SI{222}{\volt} & 744 / \textbf{\SI{723}{\volt}} & $>\SI{200}{\volt}$ \checkmark\\
$x_{\mathrm{PI}}$ & \SI{0.40}{\micro\meter} & \SI{3.07}{\micro\meter} & $>\SI{1.67}{\micro\meter}$ \checkmark\\
$a_{0,\mathrm{max}}(\omega_{r})$, $Q$ & \SI{2.7e4}{\meter\per\second\squared}, 77 & \textbf{\SI{2.8e4}{\meter\per\second\squared}}, 81 & $>\SI{2e4}{\meter\per\second\squared}$ \checkmark\\
peak $|\Delta\phi/\VAC|_{f\to0}$ (achievable) & \SI{8.7}{\milli\radian\per\volt} (222 V) & \textbf{\SI{73.9}{\milli\radian\per\volt}} (627 V) & maximized\\
\bottomrule
\end{tabular}
\end{table}

The achievable sensitivity improves by $8.5\times$ over the old design's best
stable operating point, and by $183\times$ over the \SI{10}{\volt} operating
point quoted in Part~1. One practical remark: \SI{627}{\volt} across a
\SI{0.9}{\micro\meter} gap is a $7\times10^{8}$ V/m field -- acceptable for
the implied vacuum/cryogenic quantum application, but it requires vacuum
packaging (in air, breakdown would intervene first). The voltage is high
because the fixed $t=\SI{200}{\nano\meter}$ keeps the actuator capacitance
small; the thicker device layer of part~b would also lower
$\VDCMAX\propto1/\sqrt{t}$.
